\uuid{RHDl}
\titre{Éléments simples}
\auteur{hanani ; mzari}
\organisation{exo7}
\licence{CC BY-NC-SA 4.0}
\datecreate{2020-01-20}
\niveau{L1}
\module{Analyse}
\chapitre{Calcul d'intégrales}
\sousChapitre{Primitives diverses}
\difficulte{3}
\theme{Primitives, éléments simples}
\exo7qcmid{lille-2.primitives.040}

\contenu{
\begin{question}
\qclassification{127.11}
\qtags{motcle=éléments simples}
On note par $F$ une primitive sur $\Rr$ de $\displaystyle f(x)=\frac{2x+3}{x^2+2x+2}$. Parmi les assertions suivantes, cocher celles qui sont vraies :
\begin{answers}
\good{$\forall x\in \Rr$, $\displaystyle f(x)=\frac{2x+2}{x^2+2x+2}+\frac{1}{x^2+2x+2}$.}
\bad{$\displaystyle F(x)=\ln (x^2+2x+2)+\frac{1}{x^2+2x+2}+k$, $k\in \Rr$.}
\bad{$\displaystyle F(x)=\ln (x^2+2x+2)+\arctan(x^2+2x+2)+k$, $k\in \Rr$.}
\good{$\displaystyle F(x)=\ln (x^2+2x+2)+\arctan (x+1)+k$, $k\in \Rr$.}
\end{answers}
\begin{explanations}
L'égalité $2x+3=(2x+2)+1$ donne : $\displaystyle f(x)=\frac{2x+2}{x^2+2x+2}+\frac{1}{1+(x+1)^2}$. Donc, par linéarité : $F(x)=\ln (x^2+2x+2)+\arctan (x+1)+k$ pour un $k\in \Rr$.
\end{explanations}
\end{question}
}
